%
%   Copyright (c) 2026, 2026, Hopsworks and/or its affiliates.
%
%   This program is free software; you can redistribute it and/or modify
%   it under the terms of the GNU General Public License, version 2.0,
%   as published by the Free Software Foundation.
%
%   This program is designed to work with certain software (including
%   but not limited to OpenSSL) that is licensed under separate terms,
%   as designated in a particular file or component or in included license
%   documentation.  The authors of MySQL hereby grant you an additional
%   permission to link the program and your derivative works with the
%   separately licensed software that they have either included with
%   the program or referenced in the documentation.
%
%   This program is distributed in the hope that it will be useful,
%   but WITHOUT ANY WARRANTY; without even the implied warranty of
%   MERCHANTABILITY or FITNESS FOR A PARTICULAR PURPOSE.  See the
%   GNU General Public License, version 2.0, for more details.
%
%   You should have received a copy of the GNU General Public License
%   along with this program; if not, write to the Free Software
%   Foundation, Inc., 51 Franklin St, Fifth Floor, Boston, MA 02110-1301  USA
%

\documentclass[11pt,a4paper]{report}

% Packages
\usepackage[utf8]{inputenc}
\usepackage[T1]{fontenc}
\usepackage{lmodern}
\usepackage[margin=1in]{geometry}
\usepackage{hyperref}
\usepackage{xcolor}
\usepackage{listings}
\usepackage{booktabs}
\usepackage{graphicx}
\usepackage{fancyhdr}
\usepackage{titlesec}
\usepackage{tocloft}
\usepackage{enumitem}

% Colors
\definecolor{rondb-orange}{HTML}{FFA500}
\definecolor{code-bg}{HTML}{F5F5F5}
\definecolor{code-keyword}{HTML}{0000FF}
\definecolor{code-string}{HTML}{008000}
\definecolor{code-comment}{HTML}{808080}

% Hyperref setup
\hypersetup{
    colorlinks=true,
    linkcolor=rondb-orange,
    urlcolor=blue,
    citecolor=rondb-orange,
    pdftitle={RonDB CLI Documentation},
    pdfauthor={RonDB Team},
}

% Listings setup
\lstset{
    backgroundcolor=\color{code-bg},
    basicstyle=\ttfamily\small,
    breaklines=true,
    captionpos=b,
    commentstyle=\color{code-comment},
    frame=single,
    keywordstyle=\color{code-keyword}\bfseries,
    stringstyle=\color{code-string},
    showstringspaces=false,
    tabsize=2,
    xleftmargin=1em,
    framexleftmargin=1em,
}

\lstdefinelanguage{rondb}{
    keywords={READ, BATCH, FILTER, OP, SELECT, INSERT, UPDATE, DELETE, FROM, WHERE, INTO, VALUES, SET, CREATE, DROP, TABLE, DATABASE, SHOW, DESCRIBE, USE, RONSQL, MYSQL},
    morecomment=[l]{\#},
    morestring=[b]",
    morestring=[b]',
    sensitive=false,
}

\lstdefinelanguage{bash}{
    keywords={rondb, go, build},
    morecomment=[l]{\#},
    morestring=[b]",
    morestring=[b]',
}

% Header/Footer
\pagestyle{fancy}
\fancyhf{}
\fancyhead[L]{\leftmark}
\fancyhead[R]{RonDB CLI v0.1.0}
\fancyfoot[C]{\thepage}
\renewcommand{\headrulewidth}{0.4pt}
\renewcommand{\footrulewidth}{0.4pt}

% Title formatting
\titleformat{\chapter}[display]
    {\normalfont\huge\bfseries\color{rondb-orange}}
    {\chaptertitlename\ \thechapter}{20pt}{\Huge}

% Document info
\title{
    \vspace{2cm}
    {\Huge\bfseries\color{rondb-orange} RonDB CLI}\\[0.5cm]
    {\Large Technical Documentation}\\[1cm]
    {\large Version 0.1.0}\\[2cm]
    \rule{\textwidth}{1pt}\\[0.5cm]
    {\large Rondis commands. SQL queries. One database.}\\
    \rule{\textwidth}{1pt}
}
\author{RonDB Development Team}
\date{\today}

\begin{document}

% Title page
\maketitle
\thispagestyle{empty}

% Table of contents
\tableofcontents
\newpage

%=============================================================================
\chapter{Introduction}
%=============================================================================

\section{Overview}

RonDB CLI (\texttt{rondb-cli}) is a unified command-line interface for RonDB, a distributed in-memory database that combines Redis-compatible operations with full MySQL SQL support. The CLI provides seamless interaction with four distinct client interfaces through a single interactive shell:

\begin{itemize}
    \item \textbf{MySQL} -- Full SQL support for queries, schema management, and data manipulation
    \item \textbf{REST API} -- High-performance primary key lookups using a DSL that generates JSON requests to RDRS (RonDB Data REST Server)
    \item \textbf{RonSQL} -- Aggregation queries (COUNT, SUM, AVG, MIN, MAX) executed via the REST API
    \item \textbf{Rondis} -- Redis-compatible commands for key-value operations
\end{itemize}

\subsection{Version Compatibility}

The rondb-cli source code is included in the RonDB source tree starting from RonDB 25.10. However, most functionality works with older RonDB versions:

Table: RonDB Version Compatibility

\begin{tabular}{ll}
\toprule
Interface & Minimum RonDB Version \\
\midrule
MySQL & Any version \\
REST API & 22.10 \\
Rondis & 24.10 \\
RonSQL & 24.10 \\
\bottomrule
\end{tabular}

\subsection{Design Philosophy}

The RonDB CLI follows a ``3-minute win'' philosophy:
\begin{itemize}
    \item \textbf{Connect} -- Establish connections to RonDB services quickly
    \item \textbf{Diagnose} -- Identify issues and understand system state
    \item \textbf{Decide} -- Take action based on gathered information
\end{itemize}

\noindent Core principles:
\begin{enumerate}
    \item \textbf{One CLI, no ceremony} -- Simple, unified interface
    \item \textbf{Four interfaces, one shell} -- MySQL, RonSQL, REST API, and Rondis
    \item \textbf{Reliability} -- Robust error handling and connection management
    \item \textbf{Clarity} -- Clear output formatting and helpful messages
    \item \textbf{Speed} -- Optimized for low latency operations
    \item \textbf{Composability} -- Commands can be combined and scripted
\end{enumerate}

\subsection{Key Features}

\begin{itemize}
    \item Unified shell supporting Rondis (Redis), MySQL, RonSQL, and REST API
    \item Automatic command routing based on syntax detection
    \item Interactive TUI database browser
    \item Comprehensive benchmarking suite with latency statistics
    \item Multi-line command support
    \item TLS encryption support
    \item Environment variable and command-line configuration
    \item Color-coded output with timing information
\end{itemize}

%=============================================================================
\chapter{Installation and Setup}
%=============================================================================

\section{Prerequisites}

\begin{itemize}
    \item Go 1.24.0 or later
    \item Access to a RonDB cluster with:
    \begin{itemize}
        \item Rondis service (default port: 6379)
        \item MySQL service (default port: 3306)
        \item RDRS/REST API service (default port: 4406)
    \end{itemize}
\end{itemize}

\section{Building from Source}

\begin{lstlisting}[language=bash]
# Navigate to the rondb-cli directory
cd tools/rondb-cli

# Build the binary
go build -o rondb .

# Verify the build
./rondb version
\end{lstlisting}

\section{Quick Start}

\begin{lstlisting}[language=bash]
# Connect to local RonDB with defaults
./rondb

# Connect to remote host
./rondb --host myrondb.example.com

# Connect with MySQL authentication
./rondb --host myrondb.example.com --mysql-user admin -p

# Connect with TLS enabled
./rondb --host myrondb.example.com --tls
\end{lstlisting}

%=============================================================================
\chapter{Configuration}
%=============================================================================

\section{Command-Line Flags}

Table: Command-Line Flags

\begin{tabular}{lll}
\toprule
Flag & Default & Description \\
\midrule
--host & localhost & General host for all services \\
--mysql-host & (inherits) & MySQL host (overrides --host) \\
--rdrs-host & (inherits) & REST API host (overrides --host) \\
--rondis-port & 6379 & Rondis service port \\
--mysql-port & 3306 & MySQL service port \\
--rdrs-port & 4406 & REST API service port \\
--mysql-user & root & MySQL username \\
--mysql-password & (empty) & MySQL password \\
-p, --password & false & Prompt for MySQL password \\
--tls & false & Enable TLS for all connections \\
--rdrs-tls & false & Enable TLS for REST API only \\
--rdrs-api-key & (empty) & API key for REST API authentication \\
--verbose & 0 & Verbose level \\
--no-mysql & false & Disable MySQL connection \\
--no-rdrs & false & Disable REST API connection \\
--no-rondis & false & Disable Rondis connection \\
\bottomrule
\end{tabular}

\section{Environment Variables}

Environment variables provide an alternative configuration method. Command-line flags take precedence over environment variables.

Table: Environment Variables

\begin{tabular}{ll}
\toprule
Variable & Description \\
\midrule
RONDB\_HOST & RonDB host (sets all service hosts) \\
MYSQL\_HOST & MySQL host (overrides RONDB\_HOST) \\
RDRS\_HOST & REST API host (overrides RONDB\_HOST) \\
RONDB\_RONDIS\_PORT & Rondis service port \\
RONDB\_MYSQL\_PORT & MySQL service port \\
RONDB\_RDRS\_PORT & REST API service port \\
RONDB\_MYSQL\_USER & MySQL username \\
RONDB\_MYSQL\_PASSWORD & MySQL password \\
RONDB\_RDRS\_API\_KEY & REST API authentication key \\
\bottomrule
\end{tabular}

\section{Configuration Precedence}

Configuration values are resolved in the following order (highest to lowest priority):
\begin{enumerate}
    \item Command-line flags
    \item Environment variables
    \item Default values
\end{enumerate}

%=============================================================================
\chapter{Client Interfaces}
%=============================================================================

RonDB CLI supports four distinct client interfaces, each optimized for different access patterns.

\section{MySQL (SQL Interface)}

The MySQL interface provides full SQL query support. SQL commands are detected by keywords such as \texttt{SELECT}, \texttt{INSERT}, \texttt{UPDATE}, \texttt{DELETE}, \texttt{CREATE}, \texttt{DROP}, etc.

\subsection{Supported Operations}

\begin{lstlisting}[language=rondb]
# Data queries
SELECT * FROM users WHERE id = 1;
SELECT name, email FROM users LIMIT 10;

# Data modification
INSERT INTO users (name, email) VALUES ('John', 'john@example.com');
UPDATE users SET name = 'Jane' WHERE id = 1;
DELETE FROM users WHERE id = 1;

# Schema operations
CREATE TABLE users (
    id INT PRIMARY KEY AUTO_INCREMENT,
    name VARCHAR(100),
    email VARCHAR(255)
) ENGINE=NDB;
DROP TABLE users;
ALTER TABLE users ADD COLUMN phone VARCHAR(20);

# Database operations
USE mydb;
CREATE DATABASE IF NOT EXISTS mydb;
SHOW DATABASES;
SHOW TABLES;
DESCRIBE users;
EXPLAIN SELECT * FROM users;
\end{lstlisting}

\subsection{Multi-line SQL}

SQL statements can span multiple lines. The shell detects incomplete statements and provides a continuation prompt:

\begin{lstlisting}[language=rondb]
rondb> SELECT
    ...>   id,
    ...>   name,
    ...>   email
    ...> FROM users
    ...> WHERE active = 1
    ...> ORDER BY name;
\end{lstlisting}

\subsection{Forcing MySQL Routing}

Use the \texttt{MYSQL} prefix to force command routing through MySQL:

\begin{lstlisting}[language=rondb]
rondb> MYSQL SHOW PROCESSLIST;
\end{lstlisting}

\section{REST API (DSL Interface)}

The REST API interface uses a Domain-Specific Language (DSL) for primary key read operations.

\subsection{Single READ Syntax}

\begin{lstlisting}[language=rondb]
READ database.table [col1, col2, ...] FILTER pk1=value1, pk2=value2;
\end{lstlisting}

\begin{lstlisting}[language=rondb]
# Read all columns
READ mydb.users FILTER user_id=1001;

# Read specific columns
READ mydb.users id, name, email FILTER user_id=1001;

# Composite primary key
READ mydb.order_items FILTER order_id=100, item_id=1;

# String filter value
READ mydb.products FILTER sku="ABC-123";
\end{lstlisting}

\textbf{Example with request/response:}
\begin{lstlisting}[language=rondb]
rondb> READ mydb.users id, name, email FILTER user_id=1001;

Request:
{
  "filters": [
    {"column": "user_id", "value": 1001}
  ],
  "readColumns": [
    {"column": "id", "dataReturnType": "default"},
    {"column": "name", "dataReturnType": "default"},
    {"column": "email", "dataReturnType": "default"}
  ]
}

Response:
{
  "data": {
    "id": 1001,
    "name": "John Doe",
    "email": "john@example.com"
  }
}
(12.3ms)
\end{lstlisting}

\subsection{BATCH Syntax}

For multiple operations in a single request:

\begin{lstlisting}[language=rondb]
# Syntax 1: Full READ per operation
BATCH
  READ mydb.users id, name FILTER user_id=1 OP 1
  READ mydb.users id, name FILTER user_id=2 OP 2
;

# Syntax 2: Shared header with multiple READs
BATCH mydb.users: id, name
  READ FILTER user_id=1
  READ FILTER user_id=2
;
\end{lstlisting}

\begin{lstlisting}[language=rondb]
# Batch with two filter columns (composite primary key)
BATCH READ mydb.items qty, price FILTER order_id=1, item_id=1 OP i1 READ mydb.items qty, price FILTER order_id=1, item_id=2 OP i2;
\end{lstlisting}

\textbf{Example with request/response:}
\begin{lstlisting}[language=rondb]
rondb> BATCH
    ...>   READ mydb.users id, name FILTER user_id=1 OP 1
    ...>   READ mydb.users id, name FILTER user_id=2 OP 2
    ...> ;

Request:
{
  "operations": [
    {
      "method": "POST",
      "relative-url": "mydb/users/pk-read",
      "body": {
        "filters": [{"column": "user_id", "value": 1}],
        "readColumns": [
          {"column": "id", "dataReturnType": "default"},
          {"column": "name", "dataReturnType": "default"}
        ],
        "operationId": "1"
      }
    },
    {
      "method": "POST",
      "relative-url": "mydb/users/pk-read",
      "body": {
        "filters": [{"column": "user_id", "value": 2}],
        "readColumns": [
          {"column": "id", "dataReturnType": "default"},
          {"column": "name", "dataReturnType": "default"}
        ],
        "operationId": "2"
      }
    }
  ]
}

Response:
{
  "result": [
    {"operationId": "1", "data": {"id": 1, "name": "Alice"}},
    {"operationId": "2", "data": {"id": 2, "name": "Bob"}}
  ]
}
(15.7ms)
\end{lstlisting}

\section{RonSQL (REST API SQL)}

RonSQL executes aggregation queries through the RDRS REST API, providing high-performance SQL access without MySQL overhead. Currently, RonSQL supports aggregation queries on a single table only.

\subsection{Supported Queries}

RonSQL supports aggregation functions on a single table:

\begin{lstlisting}[language=rondb]
# Count rows
RONSQL SELECT COUNT(*) FROM mydb.orders;

# Sum values
RONSQL SELECT SUM(amount) FROM mydb.orders;

# Average, min, max
RONSQL SELECT AVG(price), MIN(price), MAX(price) FROM mydb.products;

# Count with GROUP BY
RONSQL SELECT status, COUNT(*) FROM mydb.orders GROUP BY status;

# Aggregation with WHERE clause
RONSQL SELECT SUM(quantity) FROM mydb.order_items WHERE product_id = 100;

# Multiple aggregations
RONSQL SELECT COUNT(*), SUM(amount), AVG(amount) FROM mydb.transactions;
\end{lstlisting}

\textbf{Example:}
\begin{lstlisting}[language=rondb]
rondb> RONSQL SELECT status, COUNT(*) FROM mydb.orders GROUP BY status;

Request:
{
  "database": "mydb",
  "query": "SELECT status, COUNT(*) FROM orders GROUP BY status",
  "format": "JSON",
  "explain": "ALLOW"
}

Response:
{
  "data": [
    {"status": "pending", "COUNT(*)": 42},
    {"status": "shipped", "COUNT(*)": 156},
    {"status": "delivered", "COUNT(*)": 892}
  ]
}
(8.5ms)
\end{lstlisting}

\subsection{Commands}

\begin{lstlisting}[language=rondb]
# Set default database context
RONSQL SET DATABASE mydb;

# After setting database, table names don't need qualification
RONSQL SELECT COUNT(*) FROM orders;
RONSQL SELECT SUM(amount), AVG(amount) FROM transactions;
\end{lstlisting}

\subsection{Configuration}

RonSQL behavior can be configured using internal commands:

\begin{lstlisting}[language=rondb]
# Show/set database context
.ronsql_database [database_name]

# Show/set output format
# Options: JSON, JSON_ASCII, TEXT, TEXT_NOHEADER
.ronsql_format [format]

# Show/set explain mode
# Options: ALLOW, FORBID, REQUIRE, REMOVE, FORCE
.ronsql_explain [mode]
\end{lstlisting}

\subsection{Limitations}

\begin{itemize}
    \item Only single-table queries are supported (no JOINs)
    \item Only aggregation queries (COUNT, SUM, AVG, MIN, MAX)
    \item For general SQL queries, use the MySQL interface instead
\end{itemize}

\section{Rondis (Redis Protocol)}

Rondis provides Redis-compatible commands for key-value operations. Commands are automatically detected by keyword matching.

\subsection{Supported Commands}

\begin{lstlisting}[language=rondb]
# String operations
SET key value              # Store a string value
GET key                    # Retrieve a string value
DEL key [key ...]          # Delete one or more keys
MGET key [key ...]         # Get multiple values

# Numeric operations
INCR key                   # Increment integer value
DECR key                   # Decrement integer value
INCRBY key increment       # Increment by specified amount
DECRBY key decrement       # Decrement by specified amount

# Hash operations
HSET hash field value      # Set hash field
HGET hash field            # Get hash field value
HDEL hash field [field...] # Delete hash fields
HGETALL hash               # Get all fields and values

# Key operations
KEYS pattern               # Find keys matching pattern
EXISTS key [key ...]       # Check if keys exist
TTL key                    # Get time to live
EXPIRE key seconds         # Set expiration time

# Other
PING                       # Test connection
\end{lstlisting}

\subsection{Example Session}

\begin{lstlisting}[language=rondb]
rondb> SET user:1001 "John Doe"
OK (0.8ms)

rondb> GET user:1001
"John Doe" (0.5ms)

rondb> HSET user:1001:profile name "John" email "john@example.com"
(integer) 2 (1.2ms)

rondb> HGETALL user:1001:profile
1) "name"
2) "John"
3) "email"
4) "john@example.com"
(0.7ms)
\end{lstlisting}

%=============================================================================
\chapter{Internal Commands}
%=============================================================================

Internal commands are prefixed with a dot (\texttt{.}) and provide shell configuration and utility functions.

\section{Help and Information}

Table: Help Commands

\begin{tabular}{ll}
\toprule
Command & Description \\
\midrule
.help & Display general help information \\
.help internal & Display benchmark and internal commands \\
.help start & Display connection flags and environment variables \\
.tables & List all tables in accessible databases \\
\bottomrule
\end{tabular}

\section{Navigation and Exploration}

Table: Navigation Commands

\begin{tabular}{ll}
\toprule
Command & Description \\
\midrule
.browse & Open TUI database browser \\
.ui & Alias for .browse \\
.demo & Run interactive demonstration \\
\bottomrule
\end{tabular}

\section{Configuration Commands}

Table: Configuration Commands

\begin{tabular}{ll}
\toprule
Command & Description \\
\midrule
.client [N] & Set or display client ID for benchmarks \\
.debug [0|1] & Toggle debug mode \\
.verbose [0|1|2] & Set verbose level \\
.ronsql\_database [db] & Set RonSQL database context \\
.ronsql\_format [fmt] & Set RonSQL output format \\
.ronsql\_explain [mode] & Set RonSQL explain mode \\
\bottomrule
\end{tabular}

\section{Exit Commands}

\begin{lstlisting}[language=rondb]
quit    # Exit the shell
exit    # Exit the shell
q       # Exit the shell (shorthand)
\end{lstlisting}

%=============================================================================
\chapter{Benchmarking}
%=============================================================================

RonDB CLI includes a comprehensive benchmarking suite for performance testing.

\section{Benchmark Parameters}

All benchmark commands accept common parameters:

Table: Benchmark Parameters

\begin{tabular}{lll}
\toprule
Parameter & Default & Description \\
\midrule
T (threads) & 2 & Number of concurrent threads \\
N (ops) & 1000 & Number of operations per thread \\
R (rows) & 1 & Rows per operation \\
W (write\%) & 0 & Write percentage (0-100) \\
S (seconds) & 60 & Duration for continuous benchmarks \\
\bottomrule
\end{tabular}

\section{Rondis Benchmarks}

\subsection{Data Loading}

\begin{lstlisting}[language=rondb]
# Load test data
.load_rondis [T] [N] [R]

# Example: 4 threads, 5000 ops, 10 rows each = 200,000 keys
rondb> .load_rondis 4 5000 10
\end{lstlisting}

\subsection{Performance Testing}

\begin{lstlisting}[language=rondb]
# Run benchmark (T threads, N ops, R rows, W% writes)
.bench_rondis [T] [N] [R] [W]

# Run continuous benchmark for S seconds
.bench_rondis_cont [T] [N] [R] [W] [S]

# Example: 4 threads, 10000 ops, 1 row, 20% writes
rondb> .bench_rondis 4 10000 1 20
\end{lstlisting}

\subsection{Data Cleanup}

\begin{lstlisting}[language=rondb]
# Delete test data (must match load parameters)
.del_rondis [T] [N] [R]

rondb> .del_rondis 4 5000 10
\end{lstlisting}

\section{SQL Benchmarks}

SQL benchmarks use the \texttt{test.sql\_test} table with the following schema:

\begin{lstlisting}[language=rondb]
CREATE TABLE test.sql_test (
    user_id VARCHAR(64) NOT NULL,
    event_time BIGINT NOT NULL,
    description VARCHAR(128),
    value_int BIGINT DEFAULT 0,
    event_type BIGINT DEFAULT 0,
    PRIMARY KEY (user_id, event_time)
) ENGINE=NDB;
\end{lstlisting}

The table is automatically created by the \texttt{.load\_sql} command if it does not exist.

\subsection{Data Loading}

\begin{lstlisting}[language=rondb]
# Load test data into test.sql_test table
.load_sql [T] [N] [R]

rondb> .load_sql 4 5000 10
\end{lstlisting}

\subsection{Performance Testing}

\begin{lstlisting}[language=rondb]
# Run SELECT/UPDATE benchmark
.bench_sql [T] [N] [R] [W]

# Run continuous benchmark
.bench_sql_cont [T] [N] [R] [W] [S]

# Run scan-based benchmark (no writes)
.bench_sql_scan [T] [N] [R]

# Run continuous scan benchmark
.bench_sql_scan_cont [T] [N] [R] [S]
\end{lstlisting}

\subsection{Data Cleanup}

\begin{lstlisting}[language=rondb]
# Delete test data
.del_sql [T] [N] [R]

# Drop the test table entirely
.drop_sql
\end{lstlisting}

\section{REST API Benchmarks}

\begin{lstlisting}[language=rondb]
# Run batch pk-read benchmark
.bench_rdrs [T] [N] [R]

# Run continuous benchmark
.bench_rdrs_cont [T] [N] [R] [S]

# Note: Uses data from .load_sql
\end{lstlisting}

\section{Benchmark Output}

Benchmarks report comprehensive statistics:

\begin{lstlisting}
SQL Benchmark complete!
   Configuration: 4 threads x 10000 requests x 1 rows = 40000 total
   Reads: 32000 (80.0%), Writes: 8000 (20.0%)
   Throughput: 45230 ops/sec (12.5ms total)
   Latency: min=0.1ms avg=0.9ms max=15.2ms p95=2.1ms p99=4.5ms p99.9=12.3ms

Errors encountered (3 total):
   1. connection refused
   2. timeout after 30s
   3. duplicate key error
\end{lstlisting}

\section{Error Collection}

Benchmarks collect up to 100 unique errors and display them after completion. This helps identify:
\begin{itemize}
    \item Connection issues
    \item Timeout problems
    \item Data consistency errors
    \item Resource exhaustion
\end{itemize}

%=============================================================================
\chapter{TUI Database Browser}
%=============================================================================

The TUI (Terminal User Interface) browser provides interactive database exploration.

\section{Launching the Browser}

\begin{lstlisting}[language=rondb]
rondb> .browse
# or
rondb> .ui
\end{lstlisting}

\section{Interface Layout}

The browser displays a two-panel interface:

\begin{verbatim}
+---------------------------+--------------------------------+
|  Databases                |  Schema: mydb.users            |
|  ------------------------ |  ------------------------------ |
|  > mydb                   |  Column     Type      Null Key |
|      users                |  id         INT       NO   PRI |
|      orders               |  name       VARCHAR   YES      |
|    ndb_binlog_index       |  email      VARCHAR   YES      |
|  > test                   |  created_at DATETIME  YES      |
|      sql_test             |                                |
+---------------------------+--------------------------------+
\end{verbatim}

\section{Keyboard Navigation}

Table: TUI Keyboard Shortcuts

\begin{tabular}{ll}
\toprule
Key & Action \\
\midrule
j / Down & Move cursor down \\
k / Up & Move cursor up \\
h / Left & Collapse database \\
l / Right & Expand database \\
Enter / Space & Select table (show schema) \\
Tab & Toggle focus between panels \\
d & Toggle schema/data view \\
q / Esc & Quit browser \\
Ctrl+C & Quit browser \\
\bottomrule
\end{tabular}

\section{Views}

\subsection{Schema View}

Displays table column information:
\begin{itemize}
    \item Column name
    \item Data type
    \item Nullable status
    \item Key type (PRI, UNI, MUL)
    \item Default value
\end{itemize}

\subsection{Data View}

Displays sample rows from the selected table (limited to prevent overwhelming the display).

Toggle between views using the \texttt{d} key.

%=============================================================================
\chapter{Architecture}
%=============================================================================

\section{Project Structure}

\begin{verbatim}
rondb-cli/
├── main.go                 # Entry point
├── cmd/
│   └── root.go             # Cobra CLI commands
├── internal/
│   ├── client/
│   │   ├── mysql.go        # MySQL client wrapper
│   │   ├── rondis.go       # Rondis client wrapper
│   │   └── rest.go         # REST API client
│   ├── dsl/
│   │   └── parser.go       # DSL parser for READ/BATCH
│   ├── shell/
│   │   └── repl.go         # Interactive REPL
│   ├── tui/
│   │   └── browser.go      # Database browser
│   └── ui/
│       ├── colors.go       # Terminal styling
│       └── table.go        # Table formatting
├── docs/
│   └── PHILOSOPHY.md       # Design principles
└── go.mod                  # Dependencies
\end{verbatim}

\section{Client Architecture}

\subsection{MySQLClient}

\begin{lstlisting}[language=Go]
type MySQLClient struct {
    db *sql.DB
}

type MySQLOptions struct {
    Host     string
    Port     int
    User     string
    Password string
    TLS      bool
}
\end{lstlisting}

Key methods:
\begin{itemize}
    \item \texttt{Query(sql string)} -- Execute SELECT queries
    \item \texttt{Execute(sql string)} -- Execute INSERT/UPDATE/DELETE
    \item \texttt{ListDatabases()} -- Get database list
    \item \texttt{ListTablesInDB(database string)} -- Get tables
    \item \texttt{DescribeTable(database, table string)} -- Get schema
\end{itemize}

\subsection{RondisClient}

\begin{lstlisting}[language=Go]
type RondisClient struct {
    client *redis.Client
}

type RondisOptions struct {
    Host     string
    Port     int
    Password string
    TLS      bool
}
\end{lstlisting}

Key methods:
\begin{itemize}
    \item \texttt{Execute(args []string)} -- Execute any Redis command
    \item \texttt{Ping()} -- Test connection
\end{itemize}

\subsection{RestClient}

\begin{lstlisting}[language=Go]
type RestClient struct {
    client  *http.Client
    baseURL string
    apiKey  string
}

type RestOptions struct {
    Host   string
    Port   int
    TLS    bool
    APIKey string
}
\end{lstlisting}

Key methods:
\begin{itemize}
    \item \texttt{Post(endpoint string, body interface{})} -- Send POST request
    \item \texttt{Ping()} -- Test connection
\end{itemize}

\section{Command Routing}

Commands are routed based on keyword detection:

\begin{lstlisting}[language=Go]
func (s *Shell) execute(line string) error {
    lower := strings.ToLower(line)

    // 1. Internal commands (dot prefix)
    if strings.HasPrefix(lower, ".") {
        return s.executeInternal(line)
    }

    // 2. SQL commands (keyword detection)
    if isSQLCommand(lower) {
        return s.executeSQL(line)
    }

    // 3. RonSQL commands
    if strings.HasPrefix(lower, "ronsql ") {
        return s.executeRonSQL(line)
    }

    // 4. REST API READ commands
    if strings.HasPrefix(lower, "read ") {
        return s.executeREAD(line)
    }

    // 5. REST API BATCH commands
    if strings.HasPrefix(lower, "batch") {
        return s.executeBatch()
    }

    // 6. Default: Rondis
    return s.executeRondis(line)
}
\end{lstlisting}

\section{Latency Collection}

Thread-safe latency collection for benchmarks:

\begin{lstlisting}[language=Go]
type LatencyCollector struct {
    mu        sync.Mutex
    latencies []time.Duration
}

func (lc *LatencyCollector) Record(d time.Duration)
func (lc *LatencyCollector) GetTotalStats() (min, max, avg, p95, p99, p999 time.Duration, count int)
\end{lstlisting}

\section{Error Collection}

Thread-safe error collection (limited to 100 errors):

\begin{lstlisting}[language=Go]
const maxBenchmarkErrors = 100

type ErrorCollector struct {
    mu     sync.Mutex
    errors []string
    count  int64
}

func (ec *ErrorCollector) Record(err error)
func (ec *ErrorCollector) Count() int64
func (ec *ErrorCollector) PrintErrors()
\end{lstlisting}

%=============================================================================
\chapter{Dependencies}
%=============================================================================

\section{Core Dependencies}

Table: Go Module Dependencies

\begin{tabular}{lll}
\toprule
Package & Version & Purpose \\
\midrule
charmbracelet/bubbles & v0.21.0 & TUI components \\
charmbracelet/bubbletea & v1.3.10 & TUI framework \\
charmbracelet/lipgloss & v1.1.0 & Terminal styling \\
chzyer/readline & v1.5.1 & Interactive input \\
go-sql-driver/mysql & v1.7.1 & MySQL driver \\
redis/go-redis/v9 & v9.1.0 & Redis client \\
spf13/cobra & v1.7.0 & CLI framework \\
golang.org/x/term & v0.39.0 & Terminal control \\
olekukonko/tablewriter & v0.0.5 & Table formatting \\
\bottomrule
\end{tabular}

%=============================================================================
\chapter{Security Considerations}
%=============================================================================

\section{SQL Injection Prevention}

The MySQL client sanitizes identifiers to prevent SQL injection:

\begin{lstlisting}[language=Go]
func sanitizeIdentifier(name string) string {
    return strings.ReplaceAll(name, "`", "``")
}

// Usage in queries
query := fmt.Sprintf("SHOW TABLES FROM `%s`",
    sanitizeIdentifier(database))
\end{lstlisting}

\section{TLS Configuration}

TLS can be enabled for secure connections:

\begin{lstlisting}[language=Go]
tlsConfig := &tls.Config{
    MinVersion:         tls.VersionTLS12,
    InsecureSkipVerify: true,  // For self-signed certs
}
\end{lstlisting}

\textbf{Note:} \texttt{InsecureSkipVerify} is enabled by default for development environments. Production deployments should use proper certificate validation.

\section{API Key Authentication}

REST API requests include authentication headers when configured:

\begin{lstlisting}[language=Go]
if c.apiKey != "" {
    req.Header.Set("X-API-Key", c.apiKey)
}
\end{lstlisting}

\section{Password Handling}

Passwords are:
\begin{itemize}
    \item Never logged or displayed
    \item Read silently when using \texttt{-p} flag
    \item Can be provided via environment variable for automation
\end{itemize}

%=============================================================================
\chapter{Troubleshooting}
%=============================================================================

\section{Connection Issues}

\subsection{Checking Status}

Use the \texttt{status} subcommand to verify connections:

\begin{lstlisting}[language=bash]
./rondb status --host myrondb.example.com
\end{lstlisting}

\subsection{Common Connection Errors}

Table: Connection Troubleshooting

\begin{tabular}{ll}
\toprule
Error & Solution \\
\midrule
Connection refused & Verify service is running and port is correct \\
Timeout & Check network connectivity and firewall rules \\
Authentication failed & Verify username and password \\
TLS handshake error & Check TLS configuration and certificates \\
\bottomrule
\end{tabular}

\section{Verbose Modes}

Enable verbose output for debugging:

\begin{lstlisting}[language=bash]
# Level 1: Connection information
./rondb --verbose 1

# Level 2: Debug output
./rondb --verbose 2
\end{lstlisting}

\section{Debug Mode}

Enable debug mode within the shell:

\begin{lstlisting}[language=rondb]
rondb> .debug 1
Debug mode enabled

rondb> SET key value
[DEBUG] REQ: [SET key value]
[DEBUG] RESP: OK
OK (0.8ms)
\end{lstlisting}

%=============================================================================
\chapter{Examples}
%=============================================================================

\section{Complete Workflow Example}

\begin{lstlisting}[language=rondb]
# Connect to RonDB
$ ./rondb --host myrondb.example.com --mysql-user admin -p
Enter MySQL password: ********

Welcome to RonDB CLI v0.1.0
Type .help for help, .demo for a quick demo

# Run the demo
rondb> .demo

# Explore databases
rondb> .browse

# Write via Rondis
rondb> SET product:1001 '{"name":"Widget","price":29.99}'
OK (0.8ms)

# Read via Rondis
rondb> GET product:1001
{"name":"Widget","price":29.99} (0.5ms)

# Query via SQL
rondb> SELECT * FROM inventory.products WHERE id = 1001;
+------+--------+-------+
| id   | name   | price |
+------+--------+-------+
| 1001 | Widget | 29.99 |
+------+--------+-------+
1 row in set (2.3ms)

# High-performance read via REST API
rondb> READ inventory.products id, name, price FILTER id=1001;
{
  "data": {
    "id": 1001,
    "name": "Widget",
    "price": 29.99
  }
}
(1.1ms)

# Run benchmark
rondb> .load_rondis 4 10000 1
rondb> .bench_rondis 4 10000 1 20
rondb> .del_rondis 4 10000 1

# Exit
rondb> quit
\end{lstlisting}

\section{Batch Operations Example}

\begin{lstlisting}[language=rondb]
rondb> BATCH
    ...> READ inventory.products id, name FILTER id=1001 OP prod1
    ...> READ inventory.products id, name FILTER id=1002 OP prod2
    ...> READ inventory.products id, name FILTER id=1003 OP prod3
    ...> ;
{
  "result": [
    {"operationId": "prod1", "data": {"id": 1001, "name": "Widget"}},
    {"operationId": "prod2", "data": {"id": 1002, "name": "Gadget"}},
    {"operationId": "prod3", "data": {"id": 1003, "name": "Gizmo"}}
  ]
}
(3.2ms)
\end{lstlisting}

%=============================================================================
\appendix
\chapter{Quick Reference Card}
%=============================================================================

\section{Command Summary}

Table: Command Summary

\begin{tabular}{ll}
\toprule
Command & Description \\
\midrule
\multicolumn{2}{l}{\textbf{Information}} \\
.help & General help \\
.help internal & Benchmark commands \\
.help start & Connection options \\
.tables & List all tables \\
.browse & TUI browser \\
.demo & Quick demo \\
\midrule
\multicolumn{2}{l}{\textbf{Configuration}} \\
.client [N] & Set client ID \\
.debug [0|1] & Toggle debug \\
.verbose [0|1|2] & Set verbosity \\
\midrule
\multicolumn{2}{l}{\textbf{Rondis Benchmarks}} \\
.load\_rondis [T] [N] [R] & Load data \\
.bench\_rondis [T] [N] [R] [W] & Run benchmark \\
.bench\_rondis\_cont [...] [S] & Continuous \\
.del\_rondis [T] [N] [R] & Delete data \\
\midrule
\multicolumn{2}{l}{\textbf{SQL Benchmarks}} \\
.load\_sql [T] [N] [R] & Load data \\
.bench\_sql [T] [N] [R] [W] & Run benchmark \\
.bench\_sql\_cont [...] [S] & Continuous \\
.bench\_sql\_scan [T] [N] [R] & Scan benchmark \\
.del\_sql [T] [N] [R] & Delete data \\
.drop\_sql & Drop test table \\
\midrule
\multicolumn{2}{l}{\textbf{REST API Benchmarks}} \\
.bench\_rdrs [T] [N] [R] & Batch pk-read \\
.bench\_rdrs\_cont [...] [S] & Continuous \\
\midrule
\multicolumn{2}{l}{\textbf{Exit}} \\
quit / exit / q & Exit shell \\
\bottomrule
\end{tabular}

\section{Default Values}

\begin{itemize}
    \item Host: \texttt{localhost}
    \item Rondis port: \texttt{6379}
    \item MySQL port: \texttt{3306}
    \item REST API port: \texttt{4406}
    \item MySQL user: \texttt{root}
    \item Benchmark threads: \texttt{2}
    \item Benchmark operations: \texttt{1000}
    \item Benchmark rows/op: \texttt{1}
    \item Benchmark write\%: \texttt{0}
    \item Continuous duration: \texttt{60} seconds
\end{itemize}

%=============================================================================
\chapter{Version History}
%=============================================================================

\section{Version 0.1.0}

Initial release featuring:
\begin{itemize}
    \item Unified CLI for Rondis, MySQL, RonSQL, and REST API
    \item Interactive TUI database browser
    \item Comprehensive benchmarking suite
    \item DSL parser for REST API operations
    \item TLS support for all connections
    \item API key authentication for REST API
    \item Error collection for benchmarks
    \item Multi-line command support
    \item Color-coded output
\end{itemize}

\end{document}
